\section*{Chapitre \ref{ch:g2:air-around-us} --- L'air autour de nous}

\begin{solution}{exo:g2:air-around-us:1}
Par exemple~: le \omterm{def:g2:air-around-us:wind}{vent} qui plie les arbres ou fait bruire les
feuilles~; un cerf-volant ou un drapeau qui vole~; les nuages de
souffle un matin froid~; les bulles quand on souffle dans une
paille~; une porte qui claque dans un courant d'\omterm{def:g2:air-around-us:air}{air}.
\end{solution}

\begin{solution}{exo:g2:air-around-us:2}
De l'\omterm{def:g2:air-around-us:air}{air} est prisonnier à l'intérieur. En pressant, on essaie de
tasser l'\omterm{def:g2:air-around-us:air}{air} dans moins de place, et l'\omterm{def:g2:air-around-us:air}{air} repousse --- ce que
quelque chose de vraiment vide ne pourrait pas faire.
\end{solution}

\begin{solution}{exo:g2:air-around-us:3}
Le verre est plein d'\omterm{def:g2:air-around-us:air}{air}, et l'\omterm{def:g2:air-around-us:air}{air} prisonnier garde sa place~: deux
choses ne peuvent pas être au même endroit en même temps, l'eau ne
peut donc pas monter jusqu'au papier.
\end{solution}

\begin{solution}{exo:g2:air-around-us:4}
Il ne faut pas le pencher. Penché, l'\omterm{def:g2:air-around-us:air}{air} prisonnier s'échappe vers le
haut en grosses bulles, l'eau prend sa place --- et le papier est
trempé.
\end{solution}

\begin{solution}{exo:g2:air-around-us:5}
Chaque glouglou est un colis d'\omterm{def:g2:air-around-us:air}{air} qui s'échappe en bulle pendant
qu'une gorgée d'eau prend sa place. Ils passent chacun leur tour
parce que le goulot de la bouteille est étroit et que l'\omterm{def:g2:air-around-us:air}{air} et l'eau
ne peuvent pas le traverser en même temps --- ni partager la même
place à l'intérieur.
\end{solution}

\begin{solution}{exo:g2:air-around-us:6}
De l'\omterm{def:g2:air-around-us:air}{air} --- jusqu'au bord. «~Vide~» veut seulement dire «~il ne
reste rien à boire~»~; la bouteille est pleine de matière invisible.
\end{solution}

\begin{solution}{exo:g2:air-around-us:7}
Le \omterm{def:g2:air-around-us:wind}{vent} est de l'\omterm{def:g2:air-around-us:air}{air} en mouvement. Métiers~: gonfler les voiles,
faire voler les cerfs-volants, sécher le linge, faire tourner les
moulinets --- deux au choix.
\end{solution}

\begin{solution}{exo:g2:air-around-us:8}
L'\omterm{def:g2:air-around-us:air}{air} se précipite vers l'arrière par le goulot, et le ballon file
vers l'avant, dans l'autre sens.
\end{solution}

\begin{solution}{exo:g2:air-around-us:9}
La bouteille est déjà pleine d'\omterm{def:g2:air-around-us:air}{air}, et l'entonnoir scellé ne laisse
aucune issue à cet \omterm{def:g2:air-around-us:air}{air} --- l'eau ne peut donc pas entrer. La
solution~: desceller un peu le pourtour (ou soulever l'entonnoir un
instant) pour que l'\omterm{def:g2:air-around-us:air}{air} puisse s'échapper à côté~; l'eau descend
alors en glougloutant.
\end{solution}

\begin{solution}{exo:g2:air-around-us:10}
Oui~! Penche le verre plein d'\omterm{def:g2:air-around-us:air}{air} pour que son ouverture se place
sous celle de l'autre verre~: l'\omterm{def:g2:air-around-us:air}{air} s'en va en grosses bulles
argentées qui montent \emph{vers le haut} dans le verre plein d'eau
et se rassemblent à son sommet, en chassant l'eau vers le bas. L'\omterm{def:g2:air-around-us:air}{air}
grimpe et l'eau descend parce que l'\omterm{def:g2:air-around-us:air}{air} remonte toujours à travers
l'eau --- on verse, mais à l'envers.
\end{solution}
